\chapter{What Calculus Is Really About}

In the Preface, we saw that calculus turns on two geometric problems: finding the slope of a curve and finding the area under it. At first glance those two questions look completely unrelated. Let us trace both on a single concrete example -- a falling ball -- and watch how they turn out to be two sides of the same problem.

\section{Two questions that built a subject}

\subsection{The first question: how fast, right now?}

Imagine a ball dropped from a tall building. You have a stopwatch and a measuring
tape, and you record how far the ball has fallen at each instant. After $1$
second it has fallen about $4.9$ meters; after $2$ seconds, about $19.6$ meters;
after $3$ seconds, about $44.1$ meters. Galileo found, by rolling balls down ramps to slow time down enough to measure,
that the distance fallen grows like the \emph{square} of the time:
\[
  s(t) = 4.9\,t^2 \quad\text{(meters, with $t$ in seconds),}
\]
where $s(t)$ denotes the distance fallen as a function of the elapsed time $t$.
Now here is a question a child might ask and a physicist must answer: \emph{how
fast is the ball going at the exact instant $t = 2$ seconds?}

Speed is distance divided by time -- but at a
single \emph{instant} no time passes and no distance is covered. The
arithmetic comes out to $0/0$, which means nothing. ``Speed at an instant'' seems to be a
contradiction in terms. Yet your car's speedometer reads something definite at
every instant, and the ball really is moving at a definite rate as it passes the
second-floor window. Something genuine is being measured. Calculus exists, in
large part, to say exactly what.

\begin{intuitionbox}
You already know how to dodge the $0/0$ problem in practice: you measure speed
over a \emph{short} interval and call it the speed ``at'' the middle. Drive $60$
miles in one hour and your average speed is $60$ mph; but to know your speed
\emph{now} you would glance at how far you went in the last second, or the last
tenth of a second. The shorter the interval, the more the average over it
deserves to be called your speed \emph{right now}. Calculus is what happens when
you take that ``shorter and shorter'' seriously -- all the way to the limit.
\end{intuitionbox}

Let us actually do it. Average speed over an interval is
\[
  \text{average speed} = \frac{\text{change in distance}}{\text{change in time}}.
\]
We want the speed near $t = 2$. Pick a small extra time $h$ and look at the
interval from $t = 2$ to $t = 2 + h$ (the stretch of time starting at $2$
seconds and lasting $h$ seconds, written in standard interval notation as
$[2,\,2+h]$). The distance covered during that interval is $s(2+h) - s(2)$, so
\[
  \text{average speed over } [2,\,2+h]
  = \frac{s(2+h) - s(2)}{h}.
\]
With $s(t) = 4.9\,t^2$ this is
\[
  \frac{4.9(2+h)^2 - 4.9(2)^2}{h}
  = \frac{4.9\,(4 + 4h + h^2) - 4.9\cdot 4}{h}
  = \frac{4.9\,(4h + h^2)}{h}
  = 4.9\,(4 + h) = 19.6 + 4.9\,h.
\]
As long as $h \neq 0$ we were allowed to
cancel it, and we are left with $19.6 + 4.9\,h$. Now watch what happens as $h$
shrinks:

\begin{center}
\begin{tabular}{cc}
\toprule
interval length $h$ (s) & average speed $19.6 + 4.9h$ (m/s) \\
\midrule
$1$       & $24.5$ \\
$0.1$     & $20.09$ \\
$0.01$    & $19.649$ \\
$0.001$   & $19.6049$ \\
$0.0001$  & $19.60049$ \\
\bottomrule
\end{tabular}
\end{center}

As $h$ shrinks, the average speeds converge steadily toward a single number,
$19.6$ m/s. We never set $h = 0$ -- that would
have given the meaningless $0/0$ -- but we watched where the answer was
\emph{heading} as $h$ approached $0$. That destination, $19.6$ m/s, is what we
will mean by the instantaneous speed at $t = 2$. This act of watching where an
expression heads, without ever arriving, is the \textbf{limit} (the definite value
a varying quantity approaches as its input gets arbitrarily close to some target),
and it is the single technical idea on which all of calculus is built.

\begin{recoverybox}
\textbf{Pause and think:} Before reading ahead, test the formula $19.6 + 4.9h$.
What average speed would you measure over a window of just one microsecond
($h = 0.000001$ s)? You get $19.6000049$ m/s. The smaller the window, the more
the average speed is indistinguishable from $19.6$ m/s. We never have to set
$h = 0$ on the clock; we only have to read where the approach is heading.
\end{recoverybox}

\begin{happeningbox}
The fraction $\dfrac{s(2+h)-s(2)}{h}$ is genuinely $\dfrac{0}{0}$ at $h = 0$ --
there is no dodging that. What rescues us is that for every \emph{nonzero} $h$,
the fraction simplifies to the perfectly ordinary expression $19.6 + 4.9h$,
which has no trouble at all at $h = 0$. The instantaneous speed is not the value
of the bad fraction \emph{at} $h=0$ (it has none); it is the value the
\emph{good} expression is approaching. A limit asserts: ``no matter
how close you demand the output be to $19.6$, I can guarantee it by making $h$
small enough.'' That is a statement about an entire neighborhood of $h$'s, not
about the forbidden point itself. Holding those two things apart -- the value at
a point versus the value approached near it -- is an important distinction.
\end{happeningbox}

\subsection{The second question: how much, in total?}

Now flip the situation around. Suppose you do \emph{not} know the distance
fallen, but you do know the speed at every instant. From the table above we
might guess (correctly) that the falling ball's speed grows linearly with time:
\[
  v(t) = 9.8\,t \quad\text{(meters per second).}
\]
The new question is the reverse of the first: \emph{given the speed at every
instant, how far has the ball travelled in total between $t = 0$ and $t = 2$?}

If the speed were constant -- say a steady $20$ m/s -- this would be no harder
than a grocery bill: distance is speed times time, $20 \times 2 = 40$ meters.
But the speed is \emph{changing} every instant, from $0$ at the start to $19.6$
at the end. There is no single speed to multiply by. So again we are stuck, and
again the escape is to chop the problem into small pieces over which things
barely change.

\begin{intuitionbox}
Over a very short stretch of time the speed hardly changes, so over that short
stretch ``distance $=$ speed $\times$ time'' is \emph{almost} exactly right.
Slice the two seconds into many short intervals, on each one pretend the speed
holds steady, multiply to get the little distance covered, and add up all the
little distances. The thinner the slices, the better the pretense, and the
closer the sum creeps to the true total distance. Once more it is ``finer and
finer, then take the limit.''
\end{intuitionbox}

Now look at what those little distances are. Distance over a short interval is
(roughly) speed times the width of that interval -- but ``a height times a width''
is the \emph{area of a rectangle}. So adding up all those little distances is the
same as adding up the areas of thin rectangles drawn under the speed graph. The
total distance travelled \emph{is the area under the speed curve}. A question
about accumulated motion has transformed directly into a question about geometric area.

Let us add up the rectangles for $v(t) = 9.8\,t$ from $t = 0$ to $t = 2$, using
$n$ equal slices. Each slice has width
\[
  \Delta t = \frac{2 - 0}{n} = \frac{2}{n}
\]
(where the capital Greek letter delta, $\Delta$, is the standard mathematical shorthand
for ``change in'', so $\Delta t$ represents a small chunk of time).
The right edge of the $k$-th slice occurs at time $t_k = k\,\Delta t = \frac{2k}{n}$.
On that slice we evaluate the speed at its right edge:
\[
  \text{height} = v(t_k) = 9.8 \cdot \frac{2k}{n} = \frac{19.6\,k}{n}.
\]
The area of this single $k$-th rectangle is its height times its width:
\[
  \text{area of slice } k = \left(9.8 \cdot \frac{2k}{n}\right) \cdot \frac{2}{n}
  = \frac{39.2\,k}{n^2}.
\]
To get the total approximate area, we add all $n$ slices together. Mathematicians
write this sum using the capital Greek letter sigma, $\sum$, which stands for ``sum'':
$\sum_{k=1}^{n}$ means ``evaluate the expression for each step $k = 1, 2, \dots, n$
and add up all the results'':
\[
  \sum_{k=1}^{n} \underbrace{9.8 \cdot \frac{2k}{n}}_{\text{height}}
  \cdot \underbrace{\frac{2}{n}}_{\text{width}}
  = \sum_{k=1}^{n} \frac{39.2\,k}{n^2}.
\]
Because the factor $\frac{39.2}{n^2}$ does not depend on the index $k$, we can
pull it out of the summation:
\[
  \text{Total rectangle area} = \frac{39.2}{n^2}\sum_{k=1}^{n} k.
\]
To evaluate this sum, we use the classic identity for the sum of the first $n$ positive integers:
\[
  \sum_{k=1}^n k = 1 + 2 + \dots + n = \frac{n(n+1)}{2}.
\]
(You can see this geometrically: arrange the numbers $1$ to $n$ in a staircase of blocks,
pair it with an identical upside-down staircase to form an $n \times (n+1)$ rectangle of
total area $n(n+1)$, and take half).
Substituting this into our total area gives:
\[
  \text{Total rectangle area} = \frac{39.2}{n^2}\cdot\frac{n(n+1)}{2}
  = \frac{19.6\,n(n+1)}{n^2} = 19.6\cdot\frac{n+1}{n} = 19.6\left(1 + \frac{1}{n}\right).
\]
As $n \to \infty$, the excess fraction $\tfrac{1}{n}\to 0$, so the staircase sum
approaches exactly $19.6$ meters. And indeed, the area of the triangular region under
the straight line $v(t) = 9.8t$ from $t=0$ to $t=2$ is
$\frac{1}{2}\times\text{base}\times\text{height} = \frac{1}{2}(2)(19.6) = 19.6$.
Notice that $s(2) = 4.9\cdot 2^2 = 19.6$: the area under the speed graph reproduces
precisely the distance we started from in the first question (Figure~\ref{ch01:fig-rectangles}).

\begin{recoverybox}
\textbf{Pause and think:} In Figure~\ref{ch01:fig-rectangles}(a), why did our four
rectangles overestimate the true area ($24.5$ vs.\ $19.6$)? Because the ball is
speeding up, so evaluating the speed at the \emph{right} edge of each slice picked
its fastest moment during that slice. What would happen if we used the \emph{left}
edge instead? The rectangles would sit entirely below the curve, giving an
underestimate ($14.7$ meters for $n=4$). As $n \to \infty$, the gap between the
upper and lower staircase sums shrinks to zero, trapping the exact area $19.6$ in
between.
\end{recoverybox}

\begin{figure}[htbp]
\centering
\begin{tikzpicture}
  % Panel (a): Coarse approximation (n = 4)
  \begin{scope}[shift={(0,0)}]
    \begin{axis}[
      width=7.2cm, height=5.6cm,
      title={Panel (a): Coarse slices ($n=4$)},
      title style={font=\small\bfseries, text=sourceblue},
      xlabel={time $t$ (s)}, ylabel={speed $v$ (m/s)},
      xmin=0, xmax=2.2, ymin=0, ymax=22,
      axis lines=left, samples=100, domain=0:2,
      xtick={0,0.5,1,1.5,2}, ytick={0,9.8,19.6},
      yticklabels={$0$,$9.8$,$19.6$},
      enlargelimits=false, clip=false,
    ]
    \foreach \k in {1,...,4}{
      \pgfmathsetmacro\xr{0.5*\k}
      \pgfmathsetmacro\hgt{9.8*\xr}
      \edef\temp{\noexpand\addplot[draw=quotientorange!80!black, fill=quotientorange!30,
         fill opacity=0.7] coordinates
         {(\xr-0.5,0) (\xr,0) (\xr,\hgt) (\xr-0.5,\hgt)} \noexpand\closedcycle;}
      \temp
    }
    \addplot[curveblue, very thick, domain=0:2] {9.8*x};
    \node[sourceblue, font=\footnotesize, right] at (axis cs:0.1,17) {Area $\approx 24.5$};
    \node[quotientorange!85!black, font=\tiny, fill=white, inner sep=1pt] at (axis cs:0.45,12.5) {overshoot};
    \draw[->, >=stealth, quotientorange!85!black, thin] (axis cs:0.45,11.5) -- (axis cs:0.68,8.8);
    \end{axis}
  \end{scope}
  \begin{scope}[shift={(7.6,0)}]
    \begin{axis}[
      width=7.2cm, height=5.6cm,
      title={Panel (b): Finer slices ($n=12$)},
      title style={font=\small\bfseries, text=targetgreen},
      xlabel={time $t$ (s)}, ylabel={speed $v$ (m/s)},
      xmin=0, xmax=2.2, ymin=0, ymax=22,
      axis lines=left, samples=100, domain=0:2,
      xtick={0,0.5,1,1.5,2}, ytick={0,9.8,19.6},
      yticklabels={$0$,$9.8$,$19.6$},
      enlargelimits=false, clip=false,
    ]
    \foreach \k in {1,...,12}{
      \pgfmathsetmacro\xr{2.0/12.0*\k}
      \pgfmathsetmacro\xl{2.0/12.0*(\k-1)}
      \pgfmathsetmacro\hgt{9.8*\xr}
      \edef\temp{\noexpand\addplot[draw=areafill!80!black, fill=areafill!45,
         fill opacity=0.7] coordinates
         {(\xl,0) (\xr,0) (\xr,\hgt) (\xl,\hgt)} \noexpand\closedcycle;}
      \temp
    }
    \addplot[curveblue, very thick, domain=0:2] {9.8*x};
    \node[targetgreen, font=\footnotesize, right] at (axis cs:0.1,17) {Area $\approx 21.23 \to 19.6$};
    \end{axis}
  \end{scope}
\end{tikzpicture}
\caption{The total distance fallen between $t=0$ and $t=2$ equals the area under the speed graph $v(t)=9.8t$. Panel (a) uses 4 coarse right-endpoint rectangles, overestimating the area due to visible stair-step wedges. Panel (b) uses 12 thinner slices, where the staircase clings tightly to the linear curve. As $n \to \infty$, the staircase error vanishes completely, converging to the exact area of $19.6$ meters.}
\label{ch01:fig-rectangles}
\end{figure}
\section{The two big ideas, named}
Those two questions are central to calculus. Each has a name, and together they organize the rest of the book.
\begin{focusbox}
\textbf{The derivative} answers ``how fast, right now?'' It is the
\emph{instantaneous rate of change} of one quantity with respect to another --
and, drawn on a graph, it is the \emph{slope of the curve at a point}. The
process of finding it is \emph{differentiation}.
\textbf{The integral} answers ``how much, in total?'' It is the
\emph{accumulation} of a changing quantity -- and, drawn on a graph, it is the
\emph{area under the curve}. The process of finding it is \emph{integration}.
\end{focusbox}
The same two ideas appear everywhere, and recognizing them is essential.
The derivative is velocity (rate of change of position), acceleration (rate of
change of velocity), current (rate of flow of charge), marginal cost (rate of
change of cost as you make one more), reaction rate, population growth rate, and
the slope of any trend you have ever drawn through data. The integral is total
distance from a speed, total charge from a current, total cost from a marginal
cost, the work done by a varying force, the area of an irregular field, the mass
of a rod whose density varies, and the probability piled up under a bell curve.
Learn the two ideas once and you have learned them for every subject at once.
\subsection{The derivative is the slope of a tangent}
We computed the falling ball's instantaneous speed by shrinking a time interval.
Let us now \emph{see} that computation. Draw the graph of $s(t)=4.9\,t^2$ with
time across and distance up. Pick the point on the curve at $t = 2$. To estimate
the speed there we used the interval from $2$ to $2+h$; on the graph, that
amounts to drawing the straight line through the two curve points at $t=2$ and
$t=2+h$. A line through two points of a curve is called a \textbf{secant} line
(from the Latin \emph{secare}, ``to cut'', because it cuts across the curve),
and its slope -- rise over run -- is exactly
\[
  \text{slope of secant} = \frac{s(2+h)-s(2)}{(2+h)-2}
  = \frac{s(2+h)-s(2)}{h},
\]
the very same average-speed fraction as before. \emph{Average speed is the slope
of a secant line.} Now let $h \to 0$: the second point slides along the curve
toward the first, the secant pivots, and in the limit it becomes the line that
just grazes the curve at $t = 2$ -- the \textbf{tangent} line (from the Latin
\emph{tangere}, ``to touch''). Its slope is the instantaneous speed, $19.6$.
The derivative \emph{is} the slope of the tangent, and the tangent \emph{is} the
limit of secants (Figure~\ref{ch01:fig-secant-tangent}).

\begin{figure}[htbp]
\centering
\begin{tikzpicture}
\begin{axis}[
  width=11.5cm, height=7cm,
  xlabel={time $t$ (s)}, ylabel={distance $s$ (m)},
  xmin=0, xmax=3.4, ymin=0, ymax=48,
  axis lines=left, samples=100, domain=0:3.2,
  xtick={1,2,2.5,3}, ytick={0,19.6,30.6,44.1},
  yticklabels={$0$, $19.6$, $30.6$, $44.1$},
  enlargelimits=false, clip=false,
]
\addplot[curveblue, very thick] {4.9*x^2};
\node[curveblue, font=\small, above left] at (axis cs:2.7,36) {$s(t)=4.9\,t^2$};

% secant 1: through t=2 and t=3 (h=1), slope = 24.5
\addplot[gray!60, dashed, thick, domain=0.8:3.2] {19.6 + 24.5*(x-2)};
\node[gray!70!black, font=\footnotesize, right] at (axis cs:3.05,44.1) {secant ($h=1$, slope $24.5$)};

% secant 2: through t=2 and t=2.5 (h=0.5), slope = 4.9*(4+0.5) = 22.05
\addplot[quotientorange, dashed, thick, domain=0.8:3.0] {19.6 + 22.05*(x-2)};
\node[quotientorange!90!black, font=\footnotesize, right] at (axis cs:2.55,30.6) {secant ($h=0.5$, slope $22.05$)};

% tangent at t=2: slope 19.6
\addplot[tangentred, very thick, domain=0.7:3.2] {19.6 + 19.6*(x-2)};
\node[tangentred, font=\small\bfseries, right] at (axis cs:2.6,26) {tangent line (slope $19.6$)};

% points
\addplot[only marks, mark=*, mark size=2pt, black] coordinates {(2,19.6) (2.5,30.625) (3,44.1)};
\node[anchor=south east, font=\footnotesize] at (axis cs:2.0,19.6) {$(2,\,s(2))$};
\end{axis}
\end{tikzpicture}
\caption{Secant to tangent convergence. The dashed lines join $(2,s(2))$ to $(2+h,s(2+h))$ for $h=1$ and $h=0.5$; their slopes represent the average speed over each shrinking interval. As $h\to 0$, the pivot point slides along the parabola toward $(2,19.6)$ and the secant line approaches the solid red tangent line, whose slope is the instantaneous speed $19.6$ m/s.}
\label{ch01:fig-secant-tangent}
\end{figure}

\begin{whybox}
Why should ``slope of the tangent'' equal ``rate of change''? Because slope
\emph{is} a rate of change -- that is all slope ever was. The slope of a straight
line is rise over run: how much the output changes per unit change in the input.
A curve is not straight, so it has no single slope; but if you zoom in on one
point of a smooth curve, it looks straighter and straighter, until under enough
magnification it is indistinguishable from a line. That line is the tangent, and
its slope is the rate of change of the curve \emph{at that point}. The
limit-of-secants construction is just the precise way to find the slope of a
graph that refuses to be a single straight line. Every derivative measures the local slope revealed by zooming in on a curve.
\end{whybox}

\begin{walkbox}
Picture a curve that goes up steeply, then levels off, then dips down -- say a
hill in profile. Walk along it in your mind and watch the tangent line tilt as
you go: steep up, then flat at the crest, then tilting down the far side
(Figure~\ref{ch01:fig-hill}). Where is the tangent perfectly horizontal? What is
special about that point on the hill? You have just discovered, with no formulas,
why setting the derivative to zero finds the tops and bottoms -- the peaks and
valleys -- of a curve.
\end{walkbox}

\begin{figure}[htbp]
\centering
\begin{tikzpicture}
\begin{axis}[
  width=12cm, height=5.6cm,
  title style={font=\small},
  xlabel={position along the path}, ylabel={height},
  xmin=-0.2, xmax=4.3, ymin=-1.7, ymax=2.4,
  axis lines=left, samples=120, domain=0:4,
  xtick=\empty, ytick=\empty, enlargelimits=false, clip=false,
]
% a hill-then-valley profile: rises to a crest near x=1.5, dips to a trough near x=3.2
\addplot[curveblue, very thick] {1.6*exp(-2.2*(x-1.5)^2) - 1.0*exp(-3*(x-3.2)^2)};
% tangent at the crest x=1.5 (horizontal): f(1.5) ~ 1.5998
\addplot[tangentred, thick, dashed, domain=0.7:2.3] {1.5998};
\node[tangentred, font=\footnotesize, above=2pt] at (axis cs:1.5,1.65) {horizontal tangent (peak, slope $0$)};
% tangent at the trough x=3.2 (horizontal): f(3.2) ~ -0.9972
\addplot[tangentred, thick, dashed, domain=2.5:3.9] {-0.9972};
\node[tangentred, font=\footnotesize, below=3pt] at (axis cs:3.2,-1.02) {horizontal tangent (valley, slope $0$)};
% a tilted tangent on the rising part, at x=0.7
\addplot[targetgreen, thick, domain=0.1:1.3]
  {(1.6*exp(-2.2*(0.7-1.5)^2)) + (1.6*exp(-2.2*(0.7-1.5)^2))*(-2.2*2*(0.7-1.5))*(x-0.7)};
\node[targetgreen, font=\footnotesize, above right=1pt] at (axis cs:0.15,0.7) {steep climb (slope $> 0$)};
% a tilted tangent on the falling part, at x=2.3
\addplot[quotientorange, thick, domain=1.8:2.8] {0.303 - 1.853*(x-2.3)};
\node[quotientorange!90!black, font=\footnotesize, above right=1pt] at (axis cs:2.35,0.45) {descent (slope $< 0$)};
\addplot[only marks, mark=*, mark size=2pt, black]
  coordinates {(1.5,1.5998) (3.2,-0.9972) (0.7, 0.395) (2.3, 0.303)};
\end{axis}
\end{tikzpicture}
\caption{Walking the hill. On the climb the tangent tilts steeply upward (positive slope); at the crest and again at the bottom of the valley the tangent is perfectly horizontal (slope zero); on the descent the slope is negative. Those flat-tangent points identify the peaks and valleys -- which is why, once we can compute derivatives, solving $f'(x) = 0$ pinpoints the extremes of any smooth curve.}
\label{ch01:fig-hill}
\end{figure}

\subsection{The integral is the area under a curve}

The accumulation question turned distance into area, and that picture is exactly
the integral. Given any curve $y = f(x)$ sitting above the horizontal axis, the
integral of $f$ from $x = a$ to $x = b$ is the area of the region trapped between
the curve and the axis over that span. We approximate it, as we did for the
falling ball, by filling the region with thin rectangles: slice $[a,b]$ into
many narrow strips, on each strip build a rectangle whose height is the value of
$f$ somewhere on that strip, sum the rectangle areas, and then let the strips
become infinitely thin (Figure~\ref{ch01:fig-area}). The sum of rectangle areas is
called a \textbf{Riemann sum} (named after the German mathematician Bernhard Riemann,
who gave this rectangle-slicing method its strict mathematical foundation); the
limit it approaches is the integral.

\begin{figure}[htbp]
\centering
\begin{tikzpicture}
\begin{axis}[
  width=11cm, height=6cm,
  title style={font=\small},
  xlabel={$x$}, ylabel={$y$},
  xmin=-0.2, xmax=3.5, ymin=0, ymax=4.8,
  axis lines=left, samples=120, domain=0:3.1,
  xtick={0,1,2,3}, ytick=\empty,
  xticklabels={$0$,$1$,$2$,$b$},
  enlargelimits=false, clip=false,
]
\addplot[name path=curve, curveblue, very thick, domain=0:3.1] {1 + 0.35*x^2};
\foreach \k in {1,...,10}{
  \pgfmathsetmacro\xl{0.3*(\k-1)}
  \pgfmathsetmacro\xr{0.3*\k}
  \pgfmathsetmacro\hgt{1 + 0.35*\xr*\xr}
  \edef\temp{\noexpand\addplot[draw=areafill!75!black, fill=areafill!50,
     fill opacity=0.65] coordinates
     {(\xl,0) (\xr,0) (\xr,\hgt) (\xl,\hgt)} \noexpand\closedcycle;}
  \temp
}
\addplot[curveblue, very thick, domain=0:3.1] {1 + 0.35*x^2};
\node[curveblue, font=\small, right] at (axis cs:3.1,4.36) {$y=f(x)$};
\node[draw=curveblue!70, fill=white, rounded corners=2pt, font=\small, inner sep=3pt] at (axis cs:0.9,3.5)
  {$\mathrm{Total\ Area} = \displaystyle\int_a^b f(x)\,\dd x$};
\end{axis}
\end{tikzpicture}
\caption{The integral as area. The shaded rectangles approximate the area under
$y=f(x)$ from $x=a$ to $x=b$; as they get thinner their combined area converges
to the exact area, written $\int_a^b f(x)\,\dd x$. The elongated ``S'' is
Leibniz's sum sign, and $\dd x$ is the vanishing width of a strip.}
\label{ch01:fig-area}
\end{figure}

\begin{happeningbox}
The symbol $\displaystyle\int_a^b f(x)\,\dd x$ encapsulates this limiting process.
Leibniz wrote the integral sign $\int$ as an elongated letter ``S'' for \emph{summa}, Latin for ``sum,''
because an integral represents the limit of summing rectangle areas. Each approximating rectangle has area
$f(x)\,\Delta x$: a height $f(x)$ times a strip width $\Delta x$. In $\int_a^b f(x)\,\dd x$,
the factor $\dd x$ recalls the finite strip width $\Delta x$ before we pass to the limit $\Delta x \to 0$.
\end{happeningbox}

\section{The central difficulty: the infinitely small and the infinitely many}

Both constructions rely on a limiting process. To find the derivative we evaluate a ratio whose
numerator and denominator both approach zero. To find the integral we sum $n$ rectangles as the strip width
approaches zero and $n \to \infty$. Both calculations pass to a limit -- shrinking the step size $h \to 0$
for the derivative and letting the number of strips $n \to \infty$ for the integral -- yet both settle onto
exact finite numbers. Formulating this limiting process rigorously was the central challenge in the foundation
of calculus.

\begin{intuitionbox}
Why does this require care? The ratio for average rate of change, $\Delta y / \Delta x$, is undefined
when the change in input is zero: setting $\Delta x = 0$ gives $0/0$. Similarly, multiplying rate by time
fails when the rate is constantly changing. Calculus resolves this by working with finite, nonzero quantities
first -- calculating difference quotients for $h \neq 0$ or summing $n$ finite rectangles -- and then determining
the unique value those finite expressions approach. The limit is an exact finite number determined by the
behavior on a neighborhood, without evaluating at an undefined point.
\end{intuitionbox}

\begin{trapbox}
A common mistake is to set $h = 0$ prematurely. If you substitute $h = 0$ into
$\dfrac{s(2+h)-s(2)}{h}$ before algebraic simplification, you obtain the undefined expression $\dfrac{0}{0}$.
Similarly, treating an approximating strip as having width zero collapses all rectangles to zero area.
Both errors arise from treating a limiting process as evaluation at zero. Historically, trying to treat an increment as simultaneously nonzero and infinitely small caused endless confusion; limits avoid that trap by asking only where finite expressions head as the step shrinks.
\end{trapbox}

\subsection{How the pioneers got away with it}

\begin{historybox}
The two ideas are far older than the name ``calculus.'' Around $250$~BCE,
\emph{Archimedes} found the area under a parabola and approximated $\pi$ by
trapping a circle between inscribed and circumscribed polygons with up to $96$ sides---the
``method of exhaustion.'' He had the integral's core idea without its notation,
proving each result by contradiction to avoid speaking of infinity directly.

In the $1660$s and $1670$s, \emph{Isaac Newton} in England and \emph{Gottfried Wilhelm
Leibniz} in Germany independently synthesized rates of change and area into a unified
system. Newton viewed curves as paths traced by moving particles, calling his rates
``fluxions.'' Leibniz worked from differences between sequence terms, conceived of
instantaneous changes as infinitesimals ($\dd x$), and introduced the notation $\dv{y}{x}$
and $\int$ that we still use today. A bitter priority dispute broke out over who was
first, but both reached the result independently: Newton had the physical insights first,
while Leibniz built the superior notation.

What was missing was a coherent foundation. Newton spoke of quantities ``ultimately
vanishing''; Leibniz of infinitesimals that were nonzero when divided by and zero
when discarded. In $1734$, Bishop George \emph{Berkeley} exposed this inconsistency in
\emph{The Analyst}, arguing that a calculation cannot legitimately treat the same
increment as nonzero during division and zero immediately afterward. The formulas gave
correct answers, but the logic remained incomplete.

The gap stood open for a century. In the $1820$s, Augustin-Louis \emph{Cauchy} and later
Karl \emph{Weierstrass} replaced infinitesimals with the precise $\eps$--$\delta$
definition of a limit, using inequalities to define what ``approaches'' means without
infinitely small numbers. In $1960$, Abraham \emph{Robinson} showed using mathematical
logic that infinitesimals can also be made completely rigorous (nonstandard analysis).
\end{historybox}

\begin{aestheticbox}
Notice the order in which calculus grew up: physical and geometric intuition came first, and the formal $\eps$--$\delta$ foundation arrived more than a century later once mathematicians knew what they were trying to protect. Your own understanding will grow best in the same order -- see the picture first, then pin down why it must be true.
\end{aestheticbox}

\section{The central duality: differentiation and integration invert each other}

We have two processes. Differentiation takes a quantity and returns its rate of
change. Integration takes a rate and accumulates it into a total. Stated like
that, they sound like opposites -- and they are, in the precise sense that
\emph{each undoes the other}. This is the \textbf{Fundamental Theorem of
Calculus}, the central theorem of the subject, and the reason these two questions belong
in one book rather than two.

You have already watched it happen. We started with position $s(t) = 4.9t^2$.
Differentiating it (shrinking the time interval) gave the velocity
$v(t) = 9.8\,t$. Then, starting from that velocity, integrating it (summing
rectangles under its graph from $0$ to $2$) gave back $19.6$ meters $= s(2)$. We
went out by the derivative and came home by the integral. Accumulate a rate of
change over an interval and you recover the total change of the original
quantity. In symbols, with $F$ representing a quantity and $f = F'$ its rate of change
(where the prime mark $'$ is Lagrange's notation for the derivative, read ``$F$-prime''):
\[
  \int_a^b f(x)\,\dd x = F(b) - F(a).
\]

\begin{recoverybox}
\textbf{Pause and think:} Think of driving on a highway. Your speedometer gives your
instantaneous rate of change $v(t)$, while your odometer records your accumulated
distance $s(t)$. If you know your speedometer reading at every second of a two-hour
trip, you can reconstruct the exact change on the odometer by integrating:
$\int_0^2 v(t)\,\dd t = s(2) - s(0)$. Differentiation reads the speedometer from
the odometer; integration reconstructs the odometer from the speedometer.
\end{recoverybox}

\begin{figure}[htbp]
\centering
\begin{tikzpicture}
  % Multi-lane duality diagram (art_costs visual architecture)
  % Lane 1: Global Accumulated Quantity F(t)
  \node[draw=sourceblue!70, fill=sourceblue!12, rounded corners=3pt, thick,
        text width=10.2cm, align=center, inner sep=6pt] (toplane) at (0, 3.0) {
    \textbf{\small GLOBAL ACCUMULATION:}\quad $F(t)$ (Total Quantity)\\[0.25em]
    \footnotesize Overall net change across $[a,b]$: \quad $\Delta F = F(b) - F(a)$
  };

  % Lane 2: Local Instantaneous Rate f(t) = F'(t)
  \node[draw=targetgreen!70, fill=targetgreen!12, rounded corners=3pt, thick,
        text width=10.2cm, align=center, inner sep=6pt] (botlane) at (0, 0) {
    \textbf{\small LOCAL RATE OF CHANGE:}\quad $f(t) = F'(t)$ (Instantaneous Rate)\\[0.25em]
    \footnotesize Local response to a tiny nudge $\dd t$: \quad $\dd F = f(t)\,\dd t$
  };

  % Dedicated whitespace corridors for connecting arrows
  \draw[->, >=stealth, very thick, tangentred] (-2.2, 2.15) -- (-2.2, 0.85)
    node[midway, left=6pt, font=\footnotesize, align=right]
    {\textbf{\color{tangentred}Differentiation}\\[0.15em]\scriptsize $f(t) = \lim\limits_{\Delta t \to 0}\frac{\Delta F}{\Delta t}$};

  \draw[->, >=stealth, very thick, curveblue] (2.2, 0.85) -- (2.2, 2.15)
    node[midway, right=6pt, font=\footnotesize, align=left]
    {\textbf{\color{curveblue}Integration}\\[0.15em]\scriptsize $\int_a^b f(t)\,\dd t = F(b)-F(a)$};
\end{tikzpicture}
\caption{The Fundamental Theorem loop. Differentiation evaluates the local instantaneous rate of change $f(t) = F'(t)$ as a limit of difference quotients. Integration reverses the operation, accumulating the contributions $f(t)\,\dd t$ to reconstruct the total net change $F(b)-F(a)$. Differentiation and integration are exact inverse operations.}
\label{ch01:fig-duality}
\end{figure}

\begin{whybox}
Why must accumulating a rate give back the total change? Suppose $F(t)$ is your bank balance and $f(t) = F'(t)$
is the rate money flows in and out at each instant. Over a tiny slice of time
$\dd t$, the balance changes by (rate)$\times$(time) $= f(t)\,\dd t$. Add up all
those tiny changes from $a$ to $b$ -- that is the integral $\int_a^b f\,\dd t$ --
and of course you get the \emph{net} change in the balance, $F(b)-F(a)$, because
the total change is the sum of all the little changes along the way.
Differentiation decomposes a total change into local instantaneous rates; integration
reaccumulates those rates across an interval into a total change. The Fundamental Theorem of
Calculus proves that these two operations invert each other.
\end{whybox}

\begin{aestheticbox}
Tangent slopes and areas under curves look like completely different problems: one zooms in on the steepness at a single point, while the other adds up area across an entire interval. Historically they grew out of separate questions -- motion and astronomy on one side, measuring curved shapes on the other. Discovering that each undoes the other is what turned two collections of tricks into a single subject: instead of summing infinitely many rectangles by hand, we can compute areas by running differentiation backward.
\end{aestheticbox}

\begin{connectionbox}
You already know pairs of operations that undo each other: addition and
subtraction, multiplication and division, squaring and taking a square root.
Calculus adds differentiation and integration to that list, acting not on
single numbers but on entire \emph{functions} and linking local behavior (slope at a
point) with global behavior (area over an interval). Much of higher mathematics,
from differential equations to the divergence theorem of vector calculus and
differential forms, extends this same local-to-global relationship.
\end{connectionbox}

\section{Overview of the book}

We will spend the early chapters making sure we know
what a \emph{function} is, and what the special characters of calculus -- the
number $e$, the exponential $e^x$, the trigonometric functions -- actually
\emph{mean}, beyond merely manipulating formulas. Then we will build the
\emph{limit} carefully, because everything rests on it, and we will see that the
$\eps$--$\delta$ definition provides the precise
version of ``heads toward.'' With limits in hand we will construct the
derivative, learn to compute it, and understand \emph{why} each rule (the
product rule, the chain rule, and the rest) is forced upon us by the meaning of
slope. We will do the same for the integral, prove the Fundamental Theorem that
binds the two, and finally apply these tools to the equations that
govern change itself -- differential equations.

\begin{focusbox}
\emph{First}, a picture before a formula,
always: if you cannot see it, you do not yet understand it. \emph{Second}, every
rule comes with a reason -- we will never ask you to memorize a procedure whose
mechanism we have not exposed. \emph{Third}, whenever $0/0$ or an infinite sum appears, we resolve it honestly with a limit -- watching where finite quantities head rather than dividing by zero.
\end{focusbox}

\begin{labbox}
Before you turn the page, do the falling-ball computation yourself, but at a
different instant. Take $s(t) = 4.9\,t^2$ and find the instantaneous speed at
$t = 3$ by forming $\dfrac{s(3+h)-s(3)}{h}$, simplifying until the $h$ in the
denominator cancels, and reading off the limit as $h\to 0$. You should get
$29.4$ m/s. Notice that in each case the speed came out as $9.8$ times the time
-- you have just \emph{derived}, by hand, that the derivative of $4.9\,t^2$ is
$9.8\,t$.
\end{labbox}

To move from falling pebbles to more general curves, we need a precise vocabulary
for quantities that change together. In Chapter~2, we build that language: functions.

\section*{Exercises}
\addcontentsline{toc}{section}{Exercises}

\begin{exercise}[Mental check: cruise control]
If a car travels on cruise control at a steady $60$ miles per hour, its position is
$s(t) = 60t$ miles after $t$ hours. Without using any algebra, what is its average
speed over the interval $[2,\,2+h]$ for any $h > 0$? What is its instantaneous speed
at $t = 2$? Explain why, for a straight line, every secant line is already the tangent line.
\end{exercise}

\begin{exercise}[Mental check: zero slope]
A parked car sits at position $s(t) = 0$ for all time $t$. Write down the difference
quotient $\dfrac{s(2+h)-s(2)}{h}$. What is its value for any $h \neq 0$? What is the
limit as $h \to 0$? What is the slope of the horizontal tangent line?
\end{exercise}

\begin{exercise}[Average to instantaneous]
A pebble dropped from rest falls $s(t) = 4.9\,t^2$ meters in $t$ seconds.
Compute the average speed over each interval $[1,\,1+h]$ for $h = 1,\ 0.1,\
0.01,\ 0.001$. Then simplify $\dfrac{s(1+h)-s(1)}{h}$ algebraically and take the
limit as $h \to 0$. What is the instantaneous speed at $t=1$? Confirm it equals
$9.8\cdot 1$.
\end{exercise}

\begin{exercise}[A different curve]
Let $f(x) = x^2$. Form the secant slope $\dfrac{f(2+h)-f(2)}{h}$, simplify so the
denominator's $h$ cancels, and find the limit as $h\to 0$. This is the slope of
the tangent to $y=x^2$ at $x=2$. Repeat at a general point $x=a$ and show the
tangent slope is $2a$.
\end{exercise}

\begin{exercise}[Area by rectangles]
Estimate the area under $y = x$ from $x=0$ to $x=1$ by summing $n$ rectangles of
equal width $1/n$, using right-endpoint heights. Show the sum equals
$\dfrac{1}{n^2}\cdot\dfrac{n(n+1)}{2} = \dfrac{n+1}{2n}$, and take $n\to\infty$.
Check your limit against the area of the triangle from elementary geometry ($\frac{1}{2}\times 1 \times 1$).
\end{exercise}

\begin{exercise}[Conceptual: instant versus interval]
Explain in a few sentences, to a curious friend who has not taken calculus, why
``speed at a single instant'' is not simply ``distance divided by time'' at that
instant, and how the idea of a limit rescues the notion. Avoid all symbols.
\end{exercise}

\begin{exercise}[Reading the notation]
In the expression $\displaystyle\int_0^3 f(x)\,\dd x$, identify in words what each
piece refers to in the rectangle picture: the sign $\int$, the limits $0$ and
$3$, the height $f(x)$, and the width $\dd x$. Why is the integral sign shaped
like an ``S''?
\end{exercise}

\begin{exercise}[The inverse relationship, by hand]
The velocity of a cart is $v(t) = 6t$ (meters per second). Using $n$ equal
rectangles, show the area under $v$ from $0$ to $2$ approaches $12$ meters as
$n\to\infty$. Separately, guess a position function $s(t)$ whose instantaneous
rate of change is $6t$ and check that $s(2) - s(0) = 12$. You have just verified
one instance of the Fundamental Theorem.
\end{exercise}

\begin{exercise}[Where the tangent is flat]
Sketch (by hand or in your head) the graph of $y = x^2 - 4x + 3$. Using the
result that the tangent slope at $x=a$ for $y=x^2$ is $2a$, argue that the slope
of this parabola at $x=a$ is $2a - 4$. At what $x$ is the tangent horizontal?
Relate your answer to the lowest point of the parabola, and explain the
connection to the ``walk along the hill'' picture (Figure~\ref{ch01:fig-hill})
from this chapter.
\end{exercise}

\begin{exercise}[Conceptual: taming infinity]
In 1734, Bishop George Berkeley published a famous critique of early calculus (\emph{The Analyst}). In your
own words, explain what was logically uncomfortable about treating $h$ as nonzero
(to divide by it) and then as zero (to finish the calculation), and describe how
the limit avoids ever doing both to the same quantity at the same time.
\end{exercise}

\begin{exercise}[A proof in miniature]
Let $f(x) = c$ be a constant function. Using the secant-slope definition
$\dfrac{f(x+h)-f(x)}{h}$, prove carefully that the tangent slope is $0$ at every
point. Then explain geometrically why this had to be so, and physically why ``the
rate of change of something that never changes is zero'' is the only sensible
answer.
\end{exercise}

\begin{exercise}[Speed from a table]
A runner's position is recorded as $0,\ 4,\ 16,\ 36,\ 64$ meters at times
$t = 0,1,2,3,4$ seconds. Compute the average speed over each one-second interval.
Are these average speeds, or instantaneous ones? What extra information (or
limiting process) would you need to claim the runner's exact speed at $t = 2$?
Discuss why a finite table can never quite pin down an instantaneous rate.
\end{exercise}
